\newpage
\section{Missió 7: TCP o no TCP}\label{sec:mission7}

\begin{minipage}{0.10\linewidth}
% \vspace{-0.5cm}
\includegraphics[width=\linewidth]{shakespeare}
\end{minipage}
\hspace{0.03\linewidth}
\begin{minipage}{0.88\linewidth}
Anys d'èxit continuat han consolidat l'ús d'IPv4 per interconnectar xarxes dins de Mart i amb altres planetes.
Mentrestant, els equips d'enginyeria de la colònia han descobert les limitacions de l'IPv4 nu
i la necessitat d'usar almenys un dels principals protocols de transport per comunicar-se amb la Internet de la Terra.\\

\vspace{-0.25cm}
Implementar aquesta quarta capa i afegir-la a les piles de xarxa existents de cada node de la colònia és una tasca descomunal. Per això, en una primera fase només s'implementarà un dels dos protocols de transport principals.
%
La colònia et demana ajuda per comparar quantitativament aquests dos protocols i decidir quin triar.
\end{minipage}

% \vspace{0.3cm}
\noindent
\textbf{\raisebox{0.03cm}{\warning}Atenció! Aquesta missió és \raisebox{0.25em}{molt} crítica}. Pren-te-la amb calma, demana tota l'ajuda que calgui i no perdis l'ànim. Les missions anteriors es poden considerar un entrenament per a aquesta --- et recomanem molt que les acabis amb cura abans d'intentar-la.
\vspace{-0.1cm}

\subsection*{Detalls de l'informe}

L'assemblea et demana que consideris amb cura els punts següents a l'informe:

\begin{enumerate}[itemsep=0.1cm]
\item Perquè l'informe sigui autocontingut, inclou-hi la topologia de xarxa de la colònia en format CIDR quan escaigui,
i especifica valors raonables d'MTU (pots reutilitzar els de missions anteriors). Completa la descripció
especificant l'amplada de banda (velocitat de l'enllaç, en bytes per segon -- suposem que és constant).

\item Un dels dos protocols de transport candidats usa un paràmetre anomenat MSS. Indica quin i especifica el valor
més gran que permetria la comunicació entre dues màquines qualssevol de la colònia sense requerir fragmentació. Usa aquest valor a la resta de l'informe.

\item Tria un node $N_1$ d'una base connectada directament al Bus Nord. Tria un altre node $N_2$ d'una base connectada directament al Bus Sud. Per a cadascun dels dos protocols de transport principals, calcula el temps total des que $N_1$ comença a transmetre fins que $N_2$ rep el primer byte de dades del payload (càrrega útil). Per fer-ho, suposa:
  \begin{itemize}
  \item Si el protocol és orientat a connexió, les dades només es transfereixen un cop establerta la connexió.
  \item Tots els nodes i routers (encaminadors) tenen potència de càlcul infinita, és a dir, encapsular/desencapsular dades no introdueix cap retard.
  \item Les velocitats de transmissió segueixen sent finites, tal com les has definit abans a l'informe.
  \item Les taules d'encaminament són òptimes.
  \item Totes les memòries cau comencen buides per a cada protocol, i les entrades romanen vàlides un cop establertes.
  \end{itemize}
Et recomanem incloure diagrames de carrils (swimlane) per justificar la resposta.

\item Quins protocols de transport inclouen mecanismes de control de congestió? Quines condicions caldrien per permetre una transmissió contínua de $N_1$ a $N_2$ sense esperes degudes al control de congestió?

\item Indica tres casos d'ús, \textit{és a dir}, aplicacions d'usuari: el primer adequat només per a un dels protocols de transport, el segon adequat per a l'altre protocol de transport, i el tercer adequat per a tots dos.
\end{enumerate}

\vspace{-0.25cm}

\subsection*{Detalls de la retroacció}

Segueix un procés de retroacció similar al de les missions anteriors, i assegura't de tractar aquests aspectes clau:
\begin{itemize}
\item L'MSS triat és coherent amb els valors d'MTU especificats?
\item El temps de transmissió de cada paquet està ben calculat?
\item Veus algun problema amb l'encaminament, la fragmentació o l'ARP?
\item Es consideren correctament tots els elements del control de congestió?
\item Els casos d'ús proposats estan ben triats?
\end{itemize}
